% Ad Atticum 6.7 --- headnote
% ad-atticum-06-07, before the third day before the Kalends of Sextilis (~30 July) 50 BC, Tarsus

\textit{Cicero to Atticus, written at Tarsus before
the third day before the Kalends of Sextilis (before 30
July) 50 BC --- the manuscript dateline: \textit{Scr.
Tarsi ante iii K. Sext. a. 704 (50)}. A short note
from the end of the province. The succession question
of 6.5 has resolved itself for now (the next letter,
6.6, will spell out that the command has been left to
Coelius); the road home is being laid out --- the
quaestor Mescinius to be met at Laodicea for the legal
deposit of the accounts under the lex Iulia, then
Rhodes for the sake of the boys, then Athens as soon as
the etesians let him sail.}

\textit{Two private threads come into the open. The
first is family: Quintus the younger, after much
prompting and pressing on a heart that was running of
its own accord, has brought his father back to Atticus's
sister Pomponia --- the difficult marriage that has
been a worry all through book 5 and book 6. The
second is the cipher of 6.4 and 6.5 made articulate at
last: a request that Atticus look into Milo's debts and
hold him to a promise he had given. The letter closes
on Tiro, left gravely ill at Issus and now reported to
be better --- a worry Cicero will not let himself put
down, because there is no more chaste and diligent
young man in his household.}
